% ============================================================
%  part7/ch37-redshift.tex  —  Chapter 37: Redshift
% ============================================================

\chapter{Redshift as Accumulated Relaxation}
\label{ch:redshift}

\section{The provisional ansatz}

The phenomenological starting point is
\[
  1+z = \exp\!\left(\int_\gamma \mathcal{R}_{\mathrm{RSVP}}\,ds\right),
  \qquad
  \mathcal{R}_{\mathrm{RSVP}}
  = a_\Phi\Phi + a_L L - a_D D.
\]

\begin{StatusBox}{Conjectural Extension}
  This formula was motivated by the analogy between entropy
  gradient relaxation and photon frequency loss.
  It was not derived from first principles.
  This chapter states what a proper derivation would require
  and proves the structural theorem that any such derivation
  must satisfy.
\end{StatusBox}

\section{Eikonal derivation}

Assume an electromagnetic wave in the RSVP medium:
\[
  A(x,t) = a(x,t)\,e^{iS(x,t)/\varepsilon}.
\]
Insert into a modified RSVP wave equation.
At leading order in $\varepsilon$:
\[
  g_{\mathrm{eff}}^{ab}\,\partial_a S\,\partial_b S = 0.
\]

Define effective frequency $\omega = -\partial_t S$.
Along a ray:
\[
  \frac{d\omega}{ds} = -\mathcal{R}_{\mathrm{eff}}\,\omega.
\]

\begin{theorem}[Exponential redshift from linear damping]
  \label{thm:redshift-structure}
  \statusproven{}
  Any RSVP wave equation producing
  $\frac{d\omega}{ds} = -\mathcal{R}_{\mathrm{eff}}\,\omega$
  necessarily yields
  \[
    1+z = \frac{\omega_0}{\omega}
    = \exp\!\left(\int_\gamma \mathcal{R}_{\mathrm{eff}}\,ds\right).
  \]
\end{theorem}

\begin{proof}
  Solve the ODE:
  $\ln\omega(s) = \ln\omega_0 - \int_0^s \mathcal{R}_{\mathrm{eff}}\,ds'$.
  Exponentiate:
  $\omega(s)/\omega_0 = \exp(-\int_0^s \mathcal{R}_{\mathrm{eff}}\,ds')$.
  Hence $1+z = \omega_0/\omega = \exp(\int_\gamma\mathcal{R}_{\mathrm{eff}}\,ds)$.
\end{proof}

\begin{conjecture}[Redshift emergence]
  \label{conj:redshift-emergence}
  \statusconjectured{}
  A first-principles RSVP wave equation yields an effective
  relaxation rate
  \[
    \mathcal{R}_{\mathrm{eff}}
    = a_S|\nabla S| + a_\Phi|\nabla\Phi| + a_v\nabla\cdot\mathbf{v}
    + a_L L - a_D D
  \]
  with coefficients derived from the RSVP field equations.
\end{conjecture}

\begin{OpenQuestion}[Redshift derivation]
  Derive $\mathcal{R}_{\mathrm{eff}}$ by inserting an
  electromagnetic wave ansatz into the full RSVP field equations
  and taking the eikonal limit $\varepsilon \to 0$.
\end{OpenQuestion}

\section{Effective Hubble law}

\begin{definition}[Effective Hubble coefficient]
  \[
    H_{\mathrm{eff}} = c\langle\mathcal{R}_{\mathrm{eff}}\rangle,
  \]
  where $\langle\cdot\rangle$ denotes path average.
\end{definition}

\begin{proposition}[Local Hubble scaling]
  \label{prop:hubble-local}
  \statusproven{}
  For $\int_\gamma\mathcal{R}_{\mathrm{eff}}\,ds \ll 1$
  (nearby sources),
  \[
    z \approx H_{\mathrm{eff}}\,d/c.
  \]
\end{proposition}

\begin{proof}
  $e^x = 1 + x + O(x^2)$.
  For small $x = \int\mathcal{R}_{\mathrm{eff}}\,ds$:
  $z = e^x - 1 \approx x = H_{\mathrm{eff}} d/c$.
\end{proof}

\begin{ObsConseq}[Void redshift signature]
  \textbf{Mechanism:} Path-integrated RSVP relaxation.

  \textbf{Prediction:} Lines of sight through deep voids
  should show systematically lower redshift accumulation
  (path-average positivity constraint:
  $\int_\gamma\mathcal{R}_{\mathrm{eff}}\,ds \geq 0$,
  with smaller values through void interiors).

  \textbf{Comparison:} Integrated Sachs--Wolfe (ISW) effect
  measurements and void lensing surveys.

  \textbf{Status:} \statusconjectured{}
\end{ObsConseq}

\begin{WorkedExample}[First-order LCDM agreement]
  Observed: $z = 0.1$ for a source at distance $d$.

  $\Lambda$CDM: $z \approx H_0 d/c$.

  RSVP: $z = \exp(\int_\gamma \mathcal{R}_{\mathrm{eff}}\,ds) - 1
  \approx \int_\gamma\mathcal{R}_{\mathrm{eff}}\,ds$
  for small $z$.

  Both predict the same first-order value.
  Differences appear at second order and through environmental
  dependence of $\mathcal{R}_{\mathrm{eff}}$.
\end{WorkedExample}

\begin{WorkedExample}[Void vs.\ filament sightlines]
  Two galaxies at equal distance $d$.
  Sightline A crosses a deep void.
  Sightline B crosses multiple filaments.

  $\Lambda$CDM: essentially identical cosmological redshift.

  RSVP: $z_{\mathrm{void}} < z_{\mathrm{filament}}$ because
  $\langle\mathcal{R}_{\mathrm{eff}}\rangle_{\mathrm{void}}
  < \langle\mathcal{R}_{\mathrm{eff}}\rangle_{\mathrm{filament}}$.
\end{WorkedExample}

\begin{Assumptions}
  The redshift formula $1+z = \exp(\int_\gamma\mathcal{R}_{\mathrm{eff}}\,ds)$
  assumes:
  (1) High-frequency eikonal limit ($\varepsilon \to 0$).
  (2) Slowly varying RSVP background fields.
  (3) Negligible backreaction of the wave on the plenum.
  (4) Existence of an effective refractive-index description.
  Relaxing any assumption may modify the redshift law.
\end{Assumptions}

\begin{ProblemSet}
\problem{
  Show that for $\mathcal{R}_{\mathrm{eff}} = $ const, the
  RSVP redshift formula $1+z = e^{\mathcal{R}_{\mathrm{eff}} d/c}$
  reduces to $z \approx H_{\mathrm{eff}}d/c$ for $z \ll 1$.
  What is $H_{\mathrm{eff}}$ in terms of $\mathcal{R}_{\mathrm{eff}}$?
}
\problem{
  If $\mathcal{R}_{\mathrm{eff}}$ is higher in filament regions
  and lower in void regions by a factor of 2,
  estimate the redshift difference $\Delta z$ between two
  sightlines at $z \approx 0.3$ (one void, one filament).
  Compare to typical observational errors in spectroscopic surveys.
}
\problem{
  \textbf{Inverse problem.}
  A photon arrives with $z = 0.5$.
  The path length is $d = 2000$ Mpc.
  What is the average $\langle\mathcal{R}_{\mathrm{eff}}\rangle$
  along the path?
  If the path crosses $N = 10$ voids of length 200 Mpc each,
  and each void contributes $\mathcal{R}_{\mathrm{void}}$ while
  the rest contributes $\mathcal{R}_{\mathrm{avg}}$,
  write the constraint equation between
  $\mathcal{R}_{\mathrm{void}}$ and $\mathcal{R}_{\mathrm{avg}}$.
}
\end{ProblemSet}
